\newpage
\subsection{Betragsgang des Kompensierten FET-OPVs}\label{sec:1-4}
\subsubsection{Aufgabenstellung}
Der Betragsgang der Schaltung aus Abschnitt~\ref{sec:1-3}
soll am OPV mit FET-Eingangsstufe aufgenommen werden.

\subsubsection{Versuchsaufbau}
Zur Aufnahme des Betragsgangs wird die Integratorschaltung mit einer sinusförmigen Eingangsspannung variabler Frequenz aus dem Funktionsgenerator angeregt. Am Generator wurden dabei die Einstellungen "High Z" sowie eine Spitzen-Spitzen-Spannung von $U_{PP} = \SI{10}{\volt}$ gewählt. Die Eingangsspannung und die Ausgangsspannung $U_A$ werden am Oszilloskop über zwei getrennte Kanäle erfasst.

Während der Messung ist zu berücksichtigen, dass $U_A$ trotz Abgleich der Schaltung einen leichten Gleichanteil aufweisen kann. Ursache dafür ist der verbleibende Driftstrom, der den Kondensator $C$ langsam auf- oder entlädt. Deshalb wird zwischen den einzelnen Messpunkten regelmäßig neu abgeglichen, um eine Drift des Arbeitspunktes möglichst gering zu halten.

\subsubsection{Messung}

\begin{table}[h]
	\centering
	\begin{tabular}{c|c|c}
		$f \; (\mathrm{Hz})$ & $U_a \; (\mathrm{V})$ & $A \; (\mathrm{dB})$ \\
		\hline
		$1 \times 10^{0}$    & 6.88                  & -3.25                \\
		$2 \times 10^{0}$    & 3.52                  & -9.07                \\
		$5 \times 10^{0}$    & 1.44                  & -16.83               \\
		$1 \times 10^{1}$    & 0.72                  & -22.85               \\
		$2 \times 10^{1}$    & 0.36                  & -28.97               \\
		$5 \times 10^{1}$    & 0.15                  & -36.59               \\
		$1 \times 10^{2}$    & 0.08                  & -42.16               \\
	\end{tabular}
	\caption{Messungen der Ausgangsspannung für gegebene Frequenzen}\label{tab:14-ua-fuer-frequenzen}
\end{table}

\begin{figure}
	\begin{center}
		\includegraphics[width=0.55\textwidth]{14_bode_1d_Abs.png}
	\end{center}
	\caption{Betragsgang der Kompensierten Integratorschaltung}\label{fig:14-betragsgang}
\end{figure}


\subsubsection{Erkenntnisse}

Zur Phase weiß ich nicht was ich sagen soll. Die \SI{-3}{\decibel}-Grenzfrequenz ist wie man in tabelle~\ref{tab:14-ua-fuer-frequenzen}
bei etwas unter $\SI{1}{hertz}$, also schon sehr nahe DC. In Abbildung~\ref{fig:14-betragsgang} sieht man dass der betragsgang mit ca \SI{20}{\decibel} abnimmt. Das passt mit einem idealen integrator gut zusammen.
